% GENERATED FILE. Edit canonical lessons, then run npm run generate:books.
% canonical-source-hash: fnv1a64:3129f1a65069058c
% canonical-lessons: KA-C369-antu, KA-C369-kannare, KA-C369-aniriksita, KA-C369-jailu, KA-C369-kanisiko

\chapter{At Last, with My Own Eyes}
\label{ch:ka-at-last-with-my-own-eyes}
\hlchaptermodality{369}

\begin{chapteropening}
\textbf{By the end of this chapter:} \emph{I can say at last, with one's own eyes, unexpected, a jail and to appear.}

Complete the last lesson of chapter 369: I can say at last, with one's own eyes, unexpected, a jail and to appear.
\end{chapteropening}

\section[\texorpdfstring{\kn{ಅಂತೂ} (antū)}{ಅಂತೂ (antū)}]{\kn{ಅಂತೂ} (antū) — at last; anyhow}

\label{lesson:KA-C369-antu}



\begin{quote}
Before the new one: say the Kannada for a turning point, then the Kannada for to take place.
\end{quote}

\subsection*{You'll want to know: \kn{ಅಂತೂ}}
\textbf{\kn{ಅಂತೂ}} — \emph{antū} — \textquotedblleft{}at last; anyhow\textquotedblright{}.

\begin{quote}
\kn{ಅಂತೂ} \kn{ಮೂರು} \kn{ವರ್ಷದ} \kn{ನಂತರ} \kn{ಅಣ್ಣ} \kn{ಮನೆಗೆ} \kn{ಬಂದನು}.
\end{quote}

\begin{quote}
\emph{antū mūru varṣada nantara aṇṇa manege bandanu.}
\end{quote}

\begin{quote}
At last, after three years, my elder brother came home.
\end{quote}

Often with \kn{ಇಂತೂ}: \kn{ಅಂತೂ} \kn{ಇಂತೂ}, somehow or other, at long last.

\begin{cousinweb}
\textbf{\kn{ಅಂತೂ}} (\emph{antū}) is native Dravidian: it is built on the pointing base \emph{a-}, \textquotedblleft{}that, yonder\textquotedblright{} (DEDR 1), with the ending -\emph{ū}, \textquotedblleft{}even\textquotedblright{}.
\end{cousinweb}

\begin{grammarlens}[title={What you've built}]
One more word for telling what happened.
\end{grammarlens}

\subsection*{Your turn}
\begin{itemize}
\raggedright
  \item \emph{Say it:} \emph{antū}
  \item \emph{Say it:} \emph{antū}, once more
  \item \emph{Say it:} \emph{antū mūru varṣada nantara aṇṇa manege bandanu.} — At last, after three years, my elder brother came home.
  \item \emph{Recall:} say the Kannada for a turning point, then the Kannada for to take place, then say \emph{antū} again
\end{itemize}

\begin{tcolorbox}[breakable,colback=teal!4,colframe=teal!35!black,title={Before you move on}]
What does \kn{ಅಂತೂ} mean? (\textquotedblleft{}At last; anyhow\textquotedblright{}.) Say it once more.
\end{tcolorbox}
\section[\texorpdfstring{\kn{ಕಣ್ಣಾರೆ} (kaṇṇāre)}{ಕಣ್ಣಾರೆ (kaṇṇāre)}]{\kn{ಕಣ್ಣಾರೆ} (kaṇṇāre) — with one's own eyes}

\label{lesson:KA-C369-kannare}



\begin{quote}
Before the new one: say the Kannada for to take place, then the Kannada for at last.
\end{quote}

\subsection*{You'll want to know: \kn{ಕಣ್ಣಾರೆ}}
\textbf{\kn{ಕಣ್ಣಾರೆ}} — \emph{kaṇṇāre} — \textquotedblleft{}with one's own eyes\textquotedblright{}.

\begin{quote}
\kn{ನಾನು} \kn{ಆ} \kn{ಅಪಘಾತವನ್ನು} \kn{ಕಣ್ಣಾರೆ} \kn{ನೋಡಿದೆ}.
\end{quote}

\begin{quote}
\emph{nānu ā apaghātavannu kaṇṇāre nōḍide.}
\end{quote}

\begin{quote}
I saw that accident with my own eyes.
\end{quote}

\kn{ಕಣ್ಣು}, eye, and \kn{ಆರೆ}, to the full. \kn{ಮನಸಾರೆ}, with all one's heart, is built the same way.

\begin{cousinweb}
\textbf{\kn{ಕಣ್ಣಾರೆ}} (\emph{kaṇṇāre}) is a native Dravidian compound: \emph{kaṇ}, \textquotedblleft{}eye\textquotedblright{} (DEDR 1159), and \emph{āra}, \textquotedblleft{}to the full of, fully\textquotedblright{} (DEDR 368).
\end{cousinweb}

\begin{grammarlens}[title={What you've built}]
One more word for telling what happened.
\end{grammarlens}

\subsection*{Your turn}
\begin{itemize}
\raggedright
  \item \emph{Say it:} \emph{kaṇṇāre}
  \item \emph{Say it:} \emph{kaṇṇāre}, once more
  \item \emph{Say it:} \emph{nānu ā apaghātavannu kaṇṇāre nōḍide.} — I saw that accident with my own eyes.
  \item \emph{Recall:} say the Kannada for to take place, then the Kannada for at last, then say \emph{kaṇṇāre} again
\end{itemize}

\begin{tcolorbox}[breakable,colback=teal!4,colframe=teal!35!black,title={Before you move on}]
What does \kn{ಕಣ್ಣಾರೆ} mean? (\textquotedblleft{}With one's own eyes\textquotedblright{}.) Say it once more.
\end{tcolorbox}
\section[\texorpdfstring{\kn{ಅನಿರೀಕ್ಷಿತ} (anirīkṣita)}{ಅನಿರೀಕ್ಷಿತ (anirīkṣita)}]{\kn{ಅನಿರೀಕ್ಷಿತ} (anirīkṣita) — unexpected}

\label{lesson:KA-C369-aniriksita}



\begin{quote}
Before the new one: say the Kannada for at last, then the Kannada for with one's own eyes.
\end{quote}

\subsection*{You'll want to know: \kn{ಅನಿರೀಕ್ಷಿತ}}
\textbf{\kn{ಅನಿರೀಕ್ಷಿತ}} — \emph{anirīkṣita} — \textquotedblleft{}unexpected\textquotedblright{}.

\begin{quote}
\kn{ಅದು} \kn{ಒಂದು} \kn{ಅನಿರೀಕ್ಷಿತ} \kn{ಭೇಟಿ}.
\end{quote}

\begin{quote}
\emph{adu ondu anirīkṣita bhēṭi.}
\end{quote}

\begin{quote}
It was an unexpected visit.
\end{quote}

Before a noun as it is; \kn{ಅನಿರೀಕ್ಷಿತವಾಗಿ}, unexpectedly.

\begin{cousinweb}
\textbf{\kn{ಅನಿರೀಕ್ಷಿತ}} (\emph{anirīkṣita}) is borrowed from Sanskrit: the negative \emph{an-} with \emph{nirīkṣita}, \textquotedblleft{}looked for\textquotedblright{}, from \emph{nir-īkṣ}, \textquotedblleft{}to observe, to regard\textquotedblright{} (Monier-Williams).
\end{cousinweb}

\begin{grammarlens}[title={What you've built}]
One more word for telling what happened.
\end{grammarlens}

\subsection*{Your turn}
\begin{itemize}
\raggedright
  \item \emph{Say it:} \emph{anirīkṣita}
  \item \emph{Say it:} \emph{anirīkṣita}, once more
  \item \emph{Say it:} \emph{adu ondu anirīkṣita bhēṭi.} — It was an unexpected visit.
  \item \emph{Recall:} say the Kannada for at last, then the Kannada for with one's own eyes, then say \emph{anirīkṣita} again
\end{itemize}

\begin{tcolorbox}[breakable,colback=teal!4,colframe=teal!35!black,title={Before you move on}]
What does \kn{ಅನಿರೀಕ್ಷಿತ} mean? (\textquotedblleft{}Unexpected\textquotedblright{}.) Say it once more.
\end{tcolorbox}
\section[\texorpdfstring{\kn{ಜೈಲು} (jailu)}{ಜೈಲು (jailu)}]{\kn{ಜೈಲು} (jailu) — a jail}

\label{lesson:KA-C369-jailu}



\begin{quote}
Before the new one: say the Kannada for with one's own eyes, then the Kannada for unexpected.
\end{quote}

\subsection*{You'll want to know: \kn{ಜೈಲು}}
\textbf{\kn{ಜೈಲು}} — \emph{jailu} — \textquotedblleft{}a jail\textquotedblright{}.

\begin{quote}
\kn{ಅಂತೂ} \kn{ಅವನು} \kn{ಐದು} \kn{ವರ್ಷದ} \kn{ನಂತರ} \kn{ಜೈಲಿನಿಂದ} \kn{ಹೊರಗೆ} \kn{ಬಂದನು}.
\end{quote}

\begin{quote}
\emph{antū avanu aidu varṣada nantara jailininda horage bandanu.}
\end{quote}

\begin{quote}
At last, after five years, he came out of jail.
\end{quote}

Plural: \kn{ಜೈಲುಗಳು}, \emph{jailugaḷu}.

The formal word is \kn{ಕಾರಾಗೃಹ}; \kn{ಜೈಲು} \kn{ಶಿಕ್ಷೆ} is a prison sentence.

\begin{cousinweb}
\textbf{\kn{ಜೈಲು}} (\emph{jailu}) is a loan from English \emph{jail}, with the final -\emph{u} Kannada gives many loans.
\end{cousinweb}

\begin{grammarlens}[title={What you've built}]
One more word for telling what happened.
\end{grammarlens}

\subsection*{Your turn}
\begin{itemize}
\raggedright
  \item \emph{Say it:} \emph{jailu}
  \item \emph{Say it:} \emph{jailu}, once more
  \item \emph{Say it:} \emph{antū avanu aidu varṣada nantara jailininda horage bandanu.} — At last, after five years, he came out of jail.
  \item \emph{Recall:} say the Kannada for with one's own eyes, then the Kannada for unexpected, then say \emph{jailu} again
\end{itemize}

\begin{tcolorbox}[breakable,colback=teal!4,colframe=teal!35!black,title={Before you move on}]
What does \kn{ಜೈಲು} mean? (\textquotedblleft{}A jail\textquotedblright{}.) Say it once more.
\end{tcolorbox}
\section[\texorpdfstring{\kn{ಕಾಣಿಸಿಕೊ} (kāṇisiko)}{ಕಾಣಿಸಿಕೊ (kāṇisiko)}]{\kn{ಕಾಣಿಸಿಕೊ} (kāṇisiko) — to appear; to turn up}

\label{lesson:KA-C369-kanisiko}



\begin{quote}
Before the new one: say the Kannada for unexpected, then the Kannada for a jail.
\end{quote}

\subsection*{You'll want to know: \kn{ಕಾಣಿಸಿಕೊ}}
\textbf{\kn{ಕಾಣಿಸಿಕೊ}} — \emph{kāṇisiko} — \textquotedblleft{}to appear; to turn up\textquotedblright{}.

\begin{quote}
\kn{ತುಂಬಾ} \kn{ವರ್ಷದ} \kn{ನಂತರ} \kn{ಆ} \kn{ನಟ} \kn{ಟೀವಿಯಲ್ಲಿ} \kn{ಕಾಣಿಸಿಕೊಂಡನು}.
\end{quote}

\begin{quote}
\emph{tumbā varṣada nantara ā naṭa ṭīviyalli kāṇisikoṇḍanu.}
\end{quote}

\begin{quote}
After many years that actor appeared on television.
\end{quote}

Past (I …): \kn{ಕಾಣಿಸಿಕೊಂಡೆ}, \emph{kāṇisikoṇḍe}.

Of people and things showing up: \kn{ಜ್ವರ} \kn{ಕಾಣಿಸಿಕೊಂಡಿತು}, a fever appeared.

\begin{cousinweb}
\textbf{\kn{ಕಾಣಿಸಿಕೊ}} (\emph{kāṇisiko}) is a native Dravidian compound: \kn{ಕಾಣಿಸು}, \textquotedblleft{}to be seen\textquotedblright{}, from \emph{kāṇ} \textquotedblleft{}to see\textquotedblright{} (DEDR 1443), with -\emph{ko}, from \emph{koḷ}, \textquotedblleft{}to take\textquotedblright{} (DEDR 2151).
\end{cousinweb}

\begin{grammarlens}[title={What you've built}]
One more word for telling what happened.
\end{grammarlens}

\subsection*{Your turn}
\begin{itemize}
\raggedright
  \item \emph{Say it:} \emph{kāṇisiko}
  \item \emph{Say it:} \emph{kāṇisiko}, once more
  \item \emph{Say it:} \emph{tumbā varṣada nantara ā naṭa ṭīviyalli kāṇisikoṇḍanu.} — After many years that actor appeared on television.
  \item \emph{Recall:} say the Kannada for unexpected, then the Kannada for a jail, then say \emph{kāṇisiko} again
\end{itemize}

\begin{tcolorbox}[breakable,colback=teal!4,colframe=teal!35!black,title={Before you move on}]
What does \kn{ಕಾಣಿಸಿಕೊ} mean? (\textquotedblleft{}To appear; to turn up\textquotedblright{}.) Say it once more.
\end{tcolorbox}
